\documentclass[preview]{standalone}

\usepackage{amsmath}
\usepackage{amssymb}
\usepackage{stellar}
\usepackage{definitions}

\begin{document}

\id{torsion-cyclic-modules}
\genpage

\section{Annihilators and torsion}

\begin{snippetdefinition}{module-annihilator-torsion-definition}{Annihilators and torsion}
    Let \(A\) be an \integraldomain and let \(M\) be an \(A\)-\module.
    For \(v\in M\), its \emph{annihilator} is
    \[
        \moduleann(v)
        =
        \{a\in A\suchthat av=0_M\}.
    \]
    The element \(v\) is a \emph{torsion element} if
    \(\moduleann(v)\neq\{0_A\}\). The torsion subset of \(M\) is
    \[
        \moduletor(M)
        =
        \{v\in M\suchthat v\text{ is a torsion element}\}.
    \]
    The module is a \emph{torsion module} if
    \(M=\moduletor(M)\), and it is \emph{torsion-free} if
    \(\moduletor(M)=\{0_M\}\).
\end{snippetdefinition}

\begin{snippetproposition}{torsion-submodule-properties}{Properties of the torsion submodule}
    Let \(A\) be an \integraldomain and let \(M\) be an \(A\)-\module.
    Then:
    \begin{enumerate}
        \item \(\moduleann(v)\idealin A\) for every \(v\in M\);
        \item \(\moduletor(M)\submoduleleq M\);
        \item \(M/\moduletor(M)\) is torsion-free.
    \end{enumerate}
\end{snippetproposition}

\begin{snippetproof}{torsion-submodule-properties-proof}{torsion-submodule-properties}{Properties of the torsion submodule}
    First, \(0_Av=0_M\), so \(0_A\in\moduleann(v)\). If
    \(a,b\in\moduleann(v)\) and \(c\in A\), then
    \[
        (a-b)v=0_M,
        \qquad
        (ca)v=c(av)=0_M.
    \]
    Hence \(\moduleann(v)\) is an \ideal.

    Since \(\moduleann(0_M)=A\), the zero element is torsion. If
    \(v,w\in\moduletor(M)\), choose
    \(a,b\in A\difference\{0_A\}\) with \(av=0_M\) and \(bw=0_M\).
    Since \(A\) is an \integraldomain, \(ab\neq0_A\), and
    \[
        ab(v-w)=b(av)-a(bw)=0_M.
    \]
    Thus \(v-w\) is torsion. If \(c\in A\), then
    \(a(cv)=c(av)=0_M\), so \(cv\) is torsion. Therefore
    \(\moduletor(M)\submoduleleq M\).

    Finally, suppose \(v+\moduletor(M)\) is torsion in the quotient.
    Then some \(a\neq0_A\) satisfies
    \[
        av\in\moduletor(M).
    \]
    Hence some \(b\neq0_A\) satisfies \(bav=0_M\). Since
    \(ba\neq0_A\), the element \(v\) is torsion, and therefore
    \(v+\moduletor(M)=0\). The quotient is torsion-free.
\end{snippetproof}

\begin{snippetnote}{torsion-submodule-domain-hypothesis}{Why the domain hypothesis is necessary}
    If the scalar \ring has zero divisors, the product of two nonzero
    annihilating scalars may vanish. Consequently, the sum or difference
    of two torsion elements need not be torsion, so the torsion subset
    need not be a \submodule.
\end{snippetnote}

\begin{snippetproposition}{free-modules-torsion-free}{Free modules over domains are torsion-free}
    Let \(A\) be an \integraldomain. Every \freemodule \(M\) over \(A\)
    is torsion-free.
\end{snippetproposition}

\begin{snippetproof}{free-modules-torsion-free-proof}{free-modules-torsion-free}{Free modules over domains are torsion-free}
    Let \(\mathcal{B}\) be a basis of \(M\), and suppose
    \(a\in A\difference\{0_A\}\) and \(av=0_M\). Write the necessarily
    finite expansion
    \[
        v=a_1v_1+\cdots+a_rv_r
    \]
    in distinct basis elements. Then
    \[
        0_M=(aa_1)v_1+\cdots+(aa_r)v_r.
    \]
    Freeness gives \(aa_i=0_A\) for every \(i\). Since \(A\) is an
    \integraldomain and \(a\neq0_A\), every \(a_i\) vanishes. Thus
    \(v=0_M\).
\end{snippetproof}

\section{Cyclic modules}

\begin{snippetdefinition}{cyclic-module-definition}{Cyclic module}
    Let \(A\) be a \commutativering and let \(M\) be an \(A\)-\module. The
    module \(M\) is \emph{cyclic} if there exists \(v\in M\) such that
    \[
        M=Av=\{av\suchthat a\in A\}.
    \]
    Such an element \(v\) is a \emph{generator} of \(M\).
\end{snippetdefinition}

\begin{snippetexample}{polynomial-quotient-cyclic-module-example}{A cyclic polynomial quotient}
    Let \(K\) be a \field, and let
    \[
        A=K[x],
        \qquad
        I=(x^2-1),
        \qquad
        M=A/I.
    \]
    The \module \(M\) is cyclic, generated by \(1+I\), because
    \(f+I=f(1+I)\) for every \(f\in A\).

    It is also generated by \(x+I\), since
    \[
        x(x+I)=x^2+I=1+I.
    \]
\end{snippetexample}

\begin{snippetnote}{cyclic-module-versus-cyclic-additive-group}{Module and group cyclicity}
    A cyclic \(A\)-\module need not be cyclic as an additive
    \abeliangroup. For example, if \(K\) is an infinite \field, the
    additive group of
    \[
        K[x]/(x^2-1)
    \]
    cannot be generated by one element, although it is cyclic as a
    \(K[x]\)-\module.
\end{snippetnote}

\begin{snippetexample}{polynomial-quotient-torsion-module-example}{Torsion in a polynomial quotient}
    Let \(K\) be a \field, and let
    \[
        A=K[x],
        \qquad
        I=(x^2-1),
        \qquad
        M=A/I.
    \]
    For every \(v\in M\),
    \[
        (x^2-1)v=0_M.
    \]
    Hence \(\moduletor(M)=M\), so \(M\) is a \torsionmodule. In
    particular, \(M\) is not a \freemodule over the \integraldomain
    \(K[x]\).
\end{snippetexample}

\begin{snippetexample}{polynomial-quotient-annihilators-example}{Annihilators in a polynomial quotient}
    Let \(K\) be a \field with characteristic different from \(2\), and
    put
    \[
        A=K[x],
        \qquad
        I=(x^2-1),
        \qquad
        M=A/I.
    \]
    For \(v=(a+bx)+I\),
    \[
        \moduleann(v)
        =
        \begin{cases}
            A,&a=b=0,\\
            (x-1),&a=b\neq0,\\
            (x+1),&a=-b\neq0,\\
            (x^2-1),&\text{otherwise}.
        \end{cases}
    \]

    Indeed, \(I\subseteq\moduleann(v)\), and every ideal of the
    \principalidealdomain \(K[x]\) is principal. If
    \(\moduleann(v)=(g)\), then \(g\divides x^2-1\), up to a nonzero
    scalar. Moreover,
    \[
        (x-1)v=0_M
        \ifandonlyif
        x+1\divides a+bx
        \ifandonlyif
        a=b,
    \]
    while
    \[
        (x+1)v=0_M
        \ifandonlyif
        x-1\divides a+bx
        \ifandonlyif
        a=-b.
    \]
\end{snippetexample}

\begin{snippetexample}{torsion-subset-not-submodule-example}{Torsion need not form a submodule over a ring with zero divisors}
    Let \(K\) be a \field of characteristic different from \(2\), and
    let
    \[
        A=
        \left\{
            \begin{pmatrix}a&b\\b&a\end{pmatrix}
            \suchthat a,b\in K
        \right\}
        \subseteq\matrices_2(K).
    \]
    The commutative \ring \(A\) acts naturally on \(M=K^2\). For
    \(v=(x,y)^{\mathsf T}\), a matrix
    \[
        B=\begin{pmatrix}\alpha&\beta\\\beta&\alpha\end{pmatrix}
    \]
    annihilates \(v\) precisely when
    \[
        \alpha x+\beta y=0,
        \qquad
        \beta x+\alpha y=0.
    \]
    This system has a nonzero solution \((\alpha,\beta)\)
    \ifandonlyif \(x=y\) or \(x=-y\). In particular,
    \[
        \moduleann\left(\begin{pmatrix}1\\1\end{pmatrix}\right)
        =
        \linearspan_K
        \left\{\begin{pmatrix}1&-1\\-1&1\end{pmatrix}\right\},
    \]
    and
    \[
        \moduleann\left(\begin{pmatrix}1\\-1\end{pmatrix}\right)
        =
        \linearspan_K
        \left\{\begin{pmatrix}1&1\\1&1\end{pmatrix}\right\}.
    \]
    Both displayed vectors are torsion, but their sum
    \[
        \begin{pmatrix}2\\0\end{pmatrix}
    \]
    is not. Thus the torsion subset is not a \submodule. This does not
    contradict the torsion-submodule theorem because \(A\) is not an
    \integraldomain; indeed,
    \[
        \begin{pmatrix}1&1\\1&1\end{pmatrix}
        \begin{pmatrix}1&-1\\-1&1\end{pmatrix}
        =0.
    \]
\end{snippetexample}

\begin{snippetproposition}{cyclic-module-annihilator-quotient}{Cyclic modules and annihilators}
    Let \(A\) be a \commutativering, let \(M\) be a \cyclicmodule over
    \(A\), and let \(v\) be a generator. Then
    \[
        M\moduleisomorphic A/\moduleann(v).
    \]
\end{snippetproposition}

\begin{snippetproof}{cyclic-module-annihilator-quotient-proof}{cyclic-module-annihilator-quotient}{Cyclic modules and annihilators}
    Define
    \[
        \varphi\colon A\fromto M,
        \qquad
        \varphi(a)=av.
    \]
    This is an \(A\)-\modulehomomorphism and is surjective because \(v\)
    generates \(M\). Its kernel is
    \[
        \moduleker\varphi
        =
        \{a\in A\suchthat av=0_M\}
        =
        \moduleann(v).
    \]
    The first isomorphism theorem gives the result.
\end{snippetproof}

\begin{snippetproposition}{cyclic-module-generators-same-annihilator}{Generators of a cyclic module have the same annihilator}
    Let \(A\) be a \commutativering and let \(v,w\) generate the same
    \cyclicmodule over \(A\). Then
    \[
        \moduleann(v)=\moduleann(w).
    \]
\end{snippetproposition}

\begin{snippetproof}{cyclic-module-generators-same-annihilator-proof}{cyclic-module-generators-same-annihilator}{Generators of a cyclic module have the same annihilator}
    There exist \(a,b\in A\) such that \(w=av\) and \(v=bw\). If
    \(c\in\moduleann(v)\), commutativity gives
    \[
        cw=c(av)=a(cv)=0_M,
    \]
    so \(c\in\moduleann(w)\). The reverse inclusion is symmetric.
\end{snippetproof}

\begin{snippetcorollary}{cyclic-modules-same-annihilator-isomorphic}{Cyclic modules with the same annihilator}
    Let \(A\) be a \commutativering. If \cyclicmodule[cyclic modules]
    \(M,N\) over \(A\)
    have generators \(v,w\) satisfying
    \[
        \moduleann(v)=\moduleann(w),
    \]
    then \(M\moduleisomorphic N\). In particular, every cyclic module has
    the form
    \[
        A/I
    \]
    up to isomorphism, where \(I\) is the annihilator of a generator.
\end{snippetcorollary}

\begin{snippetexample}{polynomial-quotient-chinese-remainder-module-example}{Chinese remainder decomposition of a polynomial quotient}
    Let \(K\) be a \field of characteristic different from \(2\), and let
    \[
        A=K[x],
        \qquad
        I=(x^2-1),
    \]
    The ideals \(I_1=(x-1)\) and \(I_2=(x+1)\) are \comaximal, so the Chinese
    remainder theorem gives
    \[
        A/I
        \moduleisomorphic
        A/I_1\oplus A/I_2
        \moduleisomorphic
        K\oplus K.
    \]
    Under this decomposition,
    \[
        \moduleann(a+I_1,0+I_2)=I_1
        \qquad(a\neq0),
    \]
    \[
        \moduleann(0+I_1,b+I_2)=I_2
        \qquad(b\neq0),
    \]
    while, for \(a,b\neq0\),
    \[
        \moduleann(a+I_1,b+I_2)
        =
        I_1\intersection I_2
        =
        I.
    \]
\end{snippetexample}

\end{document}
